% A page that tex4ht writes with a stray end tag, which make4ht then reads
% as HTML. Each case has a character reference right after a start tag, or
% a numeric reference with a character after its semicolon. Its text must
% read as in the PDF. The addresses of the two links are guards.
\documentclass{article}
\usepackage[T1]{fontenc}
\usepackage{listings,upquote,fancyvrb,alltt,amsthm,minted}
\usepackage{hyperref}
\lstset{language=Python,basicstyle=\ttfamily}
\newtheorem{theorem}{Theorem}
\title{\& Title's text}
\author{An Author}
\begin{document}
\maketitle
\section{\& heading}
Strings: \lstinline{'Hello'}.

\begin{lstlisting}
>>> f(3)
>>> print("hi")
\end{lstlisting}

\begin{lstlisting}[language=C]
#include <stdio.h>
if (a && b < c) puts("s");
\end{lstlisting}

Inline C: \lstinline[language=C]{a<b && c>d} then $x^2$.

Verb: \verb|<a> & "b"|.

Emph: \emph{\&}. Italic: \textit{<x}. Typewriter: \texttt{<b>} and \texttt{"y"}.
Sans: \textsf{\&s}.

\begin{theorem}
\textbf{\&} in a theorem.
\end{theorem}

Footnote\footnote{\& in a footnote.} here.

Link: \href{https://example.com/?a=1&b=2}{\&here} and \href{https://example.com/it's?c=3&d=4}{quote}.

Plain: a \& b and a\&b.

\begin{Verbatim}
>>> "v" && <w>
\end{Verbatim}

\begin{verbatim}
>>> "x" && <y>
\end{verbatim}

\begin{alltt}
>>> "z" <
\end{alltt}

\begin{minted}{python}
print("Bell Circuit:")
if a < b > c: pass
\end{minted}

Mintinline: \mintinline{python}{"q" < 1}.
\ifdefined\HCode\HCode{</div>}\fi
\end{document}
